\subsection{半单模}

定义 1. —— 若一个模是一族单模的直和，则称它为半单模。\(^{(1)}\) 。

若一个多模是一族单多模的直和，则称它为半单多模（参见 I, §4, No.~4, 第 37 页，定义 7）。

A-模 M 是半单模当且仅当把它看作其位似环 \(A_{M}\) 上的模时是半单的。

例. —— 1) 化为 0 的模与单模都是半单模。

\begin{enumerate}
\def\labelenumi{\arabic{enumi})}
\setcounter{enumi}{1}
\item
  若 A 是体，则由 II, §7, No.~1, 第 292 页的定理 1，每个 A-模都是半单模。这表明一般而言半单模可以用多种方式分解为单子模的直和（但参见 VIII，第 68 页的推论 2）。
\item
  设 A 是一个主理想整环，P 是由 A 的不可约元组成的一个代表系（VII, §1, No.~3, 第 3 页）。设 M 是一个 A-模，且对每个 \(\pi \in \mathrm{P}\)，用 \(\mathrm{M}(\pi)\) 表示 \(x \in \mathrm{M}\) 中使 \(\pi x = 0\) 的那些元素组成的集合。由 VII, §2, No.~2, 第 9 页，A-模 M 是半单模当且仅当它是诸子模 \(\mathrm{M}(\pi)\) 之和；这时它是这些子模的直和。这个例子将在后文加以推广（VIII，第 65 页）。
\end{enumerate}

设 \(A_{1}\) 与 \(A_{2}\) 是交换环 K 上的代数。在 III, §4, No.~3, 第 466 页中，我们引入了代数 \(A_{1}\) 与 \(A_{2}\) 上的左双模的概念，并证明这一概念等价于环 \(A_{1} \otimes_{K} A_{2}\) 上的左模的概念。若 M 是环 \(A_{1} \otimes_{K} A_{2}\) 上的单（相应地，半单、有限生成）模，我们就称 M 是单（相应地，半单、有限生成）双模。

定理 1. —— 设 M 是一个模，它是一族单子模 \((\mathrm{S}_{i})_{i\in\mathrm{I}}\) 的（不必是直和的）和，并设 N 是 M 的一个子模。存在 I 的一个子集 J，使 M 是由 N 与诸模 \(S_{j}\)（j 取遍 J）组成的族的直和。

用 \(\mathcal{S}\) 表示由这样的子集 \(I'\)（属于 \(I\)）组成的集合：由模 \(N\) 与 \(S_i\)（\(i\) 属于 \(I'\)）组成的族之和是直和。集合 \(\mathcal{S}\) 具有有限特征：子集 \(J\)（属于 \(I\)）属于 \(\mathcal{S}\) 当且仅当 \(J\) 的每个有限子集也如此。于是集合 \(\mathcal{S}\) 有一个极大元 \(J\)（《集合论》，III, §4, No.~5, 第 171 页）。令 \(N' = N + \sum_{j \in J} S_j\) 。设 \(i\) 属于 \(I - J\)。由于 \(J\) 在 \(\mathcal{S}\) 中极大，集合 \(J \cup \{i\}\) 不属于 \(\mathcal{S}\)，从而 \(S_i \cap N' \neq 0\)。由于 \(S_i\) 是单模，我们有 \(S_i \cap N' = S_i\)。于是我们有 \(S_i \subset N'\)（对每个 \(i \in I\)），从而 \(N' = M\)。证明完毕。

推论 1. —— 每个是一族单模之和的模都是半单模。

只需把定理 1 应用于 \(\mathbf{N} = 0\) 的情形。

推论 2. —— 模 M 是半单模当且仅当 M 的每个子模都是直因子。

由定理 1，该条件是必要的。

反之，设 M 的每个子模都有余补。设 \(M'\) 是 M 的诸单子模之和，\(M''\) 是 \(M'\) 在 M 中的一个余补。假设我们有 \(M' \neq M\)，从而 \(M'' \neq 0\)。设 N 是 \(M''\) 的一个非零的单生成子模。由 VIII，第 49 页的命题 3，存在 N 的一个极大子模 P。设 Q 是 M 中与 P 互补的一个子模。则 \(N \cap Q\) 是 N 的一个子模，在 N 中与 P 互补，从而同构于 N/P（II, §1, No.~9, 第 210 页，命题 13）。于是 \(N \cap Q\) 是 \(M''\) 的一个单子模，与 \(M'\) 的定义矛盾。

于是我们有 \(\mathrm{M}' = \mathrm{M}\)，且由推论 1，模 \(\mathrm{M}\) 是半单的。

推论 3. —— 设 M 是半单模，N 是 M 的一个子模。模 N 与 M/N 都是半单的。更确切地说，若 M 是一族单模 \((\mathrm{S}_i)_{i\in \mathrm{I}}\) 的直和，则存在 I 的一个子集 J，使 M/N 同构于 \(\bigoplus_{j\in \mathrm{J}}\mathrm{S}_j\)，而 N 同构于 \(\bigoplus_{i\in \mathrm{I - J}}\mathrm{S}_i\) 。

取定理 1 中的 J。模 \(N' = \bigoplus_{j \in J} S_j\) 在 M 中与 N 互补；从而它同构于 M/N。此外，M 的子模 N 与 \(\bigoplus_{i \in I-J} S_i\) 都与 \(N'\) 互补，从而都同构于 \(M/N'\) 。

推论 4. —— 设 M 是一个半单模。则 M 是单模当且仅当 M 的自同态环 E 是体。

若 M 是单模，则由 VIII，第 47 页命题 2 的推论，E 是体。

若 E 是体，则模 M 不可分解（VIII，第 31 页，命题 4, a)）。由于它还是半单的，故它是单模。

注. —— 设 K 是一个代数闭交换域，A 是一个 K-代数。设 M 是一个半单 A-模，且是体 K 上的有限维向量空间。则 M 是单模当且仅当 A-模 M 的每个自同态都形如 \(x \mapsto \alpha x\)（\(\alpha\) 属于 K）：其必要性由 VIII，第 47 页的定理 1 得出，充分性由上面的推论 4 得出。

